
% Mathematical Formulations of Classical Mechanics: Comparative LaTeX Monograph

\documentclass[11pt]{book}

\usepackage[a4paper,margin=1in]{geometry}
\usepackage{amsmath,amssymb,amsfonts,amsthm,mathtools}
\usepackage{physics}
\usepackage{bm}
\usepackage{mathrsfs}
\usepackage{tensor}
\usepackage{graphicx}
\usepackage{hyperref}
\usepackage{longtable}
\usepackage{booktabs}
\usepackage{enumitem}
\usepackage{tikz-cd}
\usepackage{slashed}
\usepackage{cite}

\title{Mathematical Formulations of Classical Mechanics\\
A Comparative Monograph of Geometric, Variational, Statistical, and Algebraic Structures}

\author{}
\date{}

\begin{document}

\maketitle
\tableofcontents

\chapter*{Preface}

Classical mechanics is often presented as a single theory possessing several equivalent mathematical representations.

In reality, classical mechanics admits a remarkably rich hierarchy of formulations based upon fundamentally different mathematical structures:

\begin{itemize}
\item Newtonian vector dynamics,
\item variational principles,
\item symplectic geometry,
\item Hamilton--Jacobi theory,
\item contact geometry,
\item Poisson geometry,
\item kinetic theory,
\item hydrodynamics,
\item geometric mechanics,
\item Lie-group methods,
\item Koopman operator theory,
\item information geometry,
\item category-theoretic formulations,
\item integrable-system methods,
\item stochastic mechanics.
\end{itemize}

These are not merely calculational variants. Each formulation emphasizes different primitive notions:

\begin{center}
\begin{tabular}{ll}
Newtonian mechanics & forces and trajectories \\
Lagrangian mechanics & action functionals \\
Hamiltonian mechanics & symplectic flow \\
Hamilton--Jacobi theory & generating functions \\
Poisson mechanics & Lie algebras \\
Geometric mechanics & manifolds and momentum maps \\
Kinetic theory & probability densities \\
Hydrodynamics & continuum fields \\
Koopman mechanics & Hilbert-space evolution \\
Integrable systems & spectral structures
\end{tabular}
\end{center}

The purpose of this monograph is to compare these formulations systematically from the viewpoint of mathematical structure.

\part{Foundational Structures}

\chapter{Configuration Space and Kinematics}

\section{Configuration Manifolds}

The configuration space of a mechanical system is a differentiable manifold:

\[
Q.
\]

Local coordinates:

\[
q^i.
\]

Examples:

\begin{itemize}
\item particle in Euclidean space:
\[
Q=\mathbb R^3,
\]

\item rigid body:
\[
Q=SO(3),
\]

\item pendulum:
\[
Q=S^1,
\]

\item double pendulum:
\[
Q=S^1\times S^1.
\]
\end{itemize}

\section{Velocity Phase Space}

Velocities belong to tangent bundle:

\[
TQ.
\]

Coordinates:

\[
(q^i,\dot q^i).
\]

Acceleration defines second-order vector field:

\[
X:TQ\to TTQ.
\]

\chapter{Newtonian Mechanics}

\section{Newton's Laws}

Dynamics governed by:

\[
m_i \ddot q^i = F^i.
\]

For conservative systems:

\[
F^i = -\frac{\partial V}{\partial q^i}.
\]

Hence:

\[
m_i \ddot q^i
+
\frac{\partial V}{\partial q^i}
=0.
\]

\section{Geometry of Newtonian Mechanics}

Newtonian mechanics naturally lives on tangent bundle:

\[
TQ.
\]

Dynamics becomes flow of second-order differential equations.

\section{Conservation Laws}

Momentum conservation:

\[
\frac{dp}{dt}=0.
\]

Energy conservation:

\[
E=T+V.
\]

Angular momentum:

\[
L=q\times p.
\]

These emerge from symmetry.

\part{Variational Mechanics}

\chapter{Lagrangian Mechanics}

\section{Action Principle}

Lagrangian:

\[
L:TQ\to\mathbb R.
\]

Action functional:

\[
S[q]
=
\int_{t_1}^{t_2}
L(q,\dot q,t)dt.
\]

Physical trajectories satisfy:

\[
\delta S=0.
\]

\section{Euler--Lagrange Equations}

Variation gives:

\[
\frac{d}{dt}
\frac{\partial L}{\partial \dot q^i}
-
\frac{\partial L}{\partial q^i}
=0.
\]

For:

\[
L=T-V,
\]

Newton equations are recovered.

\section{Coordinate Invariance}

Lagrangian mechanics is naturally covariant under coordinate transformations.

This makes it ideal for:

\begin{itemize}
\item constrained systems,
\item curved manifolds,
\item field theories,
\item relativity.
\end{itemize}

\chapter{Noether Theory}

\section{Continuous Symmetries}

Suppose:

\[
q^i \to q^i+\epsilon X^i(q).
\]

If action invariant:

\[
\delta S=0,
\]

then conserved quantity exists.

\section{Noether Current}

Conserved quantity:

\[
J
=
\sum_i
\frac{\partial L}{\partial \dot q^i}
X^i.
\]

\section{Examples}

Time translation:

\[
\Rightarrow
\text{energy conservation}.
\]

Spatial translation:

\[
\Rightarrow
\text{momentum conservation}.
\]

Rotational symmetry:

\[
\Rightarrow
\text{angular momentum conservation}.
\]

\part{Hamiltonian and Symplectic Mechanics}

\chapter{Hamiltonian Mechanics}

\section{Cotangent Bundle}

Phase space is cotangent bundle:

\[
T^*Q.
\]

Coordinates:

\[
(q^i,p_i).
\]

Canonical one-form:

\[
\theta
=
p_i dq^i.
\]

Canonical symplectic form:

\[
\omega
=
dq^i\wedge dp_i.
\]

\section{Legendre Transform}

Momenta:

\[
p_i
=
\frac{\partial L}{\partial \dot q^i}.
\]

Hamiltonian:

\[
H(q,p)
=
p_i\dot q^i-L.
\]

\section{Hamilton Equations}

Dynamics:

\begin{align}
\dot q^i &= \frac{\partial H}{\partial p_i},\\
\dot p_i &= -\frac{\partial H}{\partial q^i}.
\end{align}

\section{Hamiltonian Vector Fields}

Hamiltonian flow satisfies:

\[
i_{X_H}\omega=dH.
\]

Thus mechanics becomes symplectic geometry.

\chapter{Poisson Geometry}

\section{Poisson Brackets}

Define:

\[
\{f,g\}
=
\omega(X_f,X_g).
\]

Canonical form:

\[
\{f,g\}
=
\sum_i
\left(
\frac{\partial f}{\partial q^i}
\frac{\partial g}{\partial p_i}
-
\frac{\partial f}{\partial p_i}
\frac{\partial g}{\partial q^i}
\right).
\]

\section{Lie Algebra Structure}

Poisson brackets satisfy:

\begin{align}
\{f,g\} &=-\{g,f\},\\
\{f,gh\} &= \{f,g\}h+g\{f,h\},\\
\{f,\{g,h\}\} &+ \text{cyclic}=0.
\end{align}

Hence observables form Poisson algebra.

\section{General Poisson Manifolds}

Poisson tensor:

\[
\Pi^{ij}.
\]

Bracket:

\[
\{f,g\}
=
\Pi^{ij}
\partial_i f
\partial_j g.
\]

This generalizes symplectic geometry.

\chapter{Hamilton--Jacobi Theory}

\section{Hamilton Principal Function}

Action:

\[
S(q,t).
\]

Hamilton--Jacobi equation:

\[
\frac{\partial S}{\partial t}
+
H\left(q,\frac{\partial S}{\partial q}\right)
=0.
\]

\section{Characteristics}

Characteristics satisfy Hamilton equations.

Hence HJ theory converts ODE mechanics into PDE theory.

\section{Canonical Transformations}

Generating function:

\[
F(q,Q,t).
\]

Transformation equations:

\begin{align}
p_i &= \frac{\partial F}{\partial q^i},\\
P_i &= -\frac{\partial F}{\partial Q^i}.
\end{align}

\section{Action-Angle Variables}

For integrable systems:

\[
(I_i,\theta_i).
\]

Dynamics simplifies to:

\begin{align}
\dot I_i &=0,\\
\dot \theta_i &= \omega_i(I).
\end{align}

\part{Geometric Mechanics}

\chapter{Symplectic Geometry}

\section{Symplectic Manifolds}

A symplectic manifold is:

\[
(M,\omega)
\]

with:

\begin{align}
d\omega &=0,\\
\omega^n &\neq 0.
\end{align}

\section{Darboux Theorem}

Locally:

\[
\omega
=
\sum_i dq^i\wedge dp_i.
\]

Thus all symplectic manifolds locally equivalent.

\section{Symplectic Flow}

Hamiltonian flow preserves symplectic form:

\[
\mathcal L_{X_H}\omega=0.
\]

Hence phase-space volume preserved.

\chapter{Momentum Maps and Reduction}

\section{Lie Group Actions}

Suppose Lie group:

\[
G
\]

acts symplectically on:

\[
(M,\omega).
\]

\section{Momentum Map}

Momentum map:

\[
\mu:M\to\mathfrak g^*.
\]

satisfies:

\[
d\langle\mu,\xi\rangle
=
i_{X_\xi}\omega.
\]

\section{Marsden--Weinstein Reduction}

Reduced phase space:

\[
M_{red}
=
\mu^{-1}(0)/G.
\]

Gauge theories naturally arise this way.

\chapter{Contact Geometry}

\section{Odd-Dimensional Geometry}

Contact manifold:

\[
(M,\alpha)
\]

with:

\[
\alpha\wedge(d\alpha)^n\neq0.
\]

\section{Dissipative Systems}

Contact Hamilton equations:

\begin{align}
\dot q^i &= \frac{\partial H}{\partial p_i},\\
\dot p_i &= -\frac{\partial H}{\partial q^i}-p_i\frac{\partial H}{\partial s},\\
\dot s &= p_i\frac{\partial H}{\partial p_i}-H.
\end{align}

This geometrizes friction and thermodynamics.

\part{Statistical and Kinetic Formulations}

\chapter{Liouville Mechanics}

\section{Phase-Space Densities}

Probability density:

\[
\rho(q,p,t).
\]

\section{Liouville Equation}

\[
\partial_t\rho
+
\{\rho,H\}
=0.
\]

Equivalent form:

\[
\frac{d\rho}{dt}=0.
\]

\section{Ergodic Theory}

Important concepts:

\begin{itemize}
\item invariant measures,
\item mixing,
\item recurrence,
\item entropy.
\end{itemize}

\chapter{BBGKY and Kinetic Theory}

\section{Distribution Functions}

Many-body systems described by:

\[
f_N(q_1,p_1,\dots,q_N,p_N).
\]

Reduced distributions:

\[
f_n.
\]

\section{BBGKY Hierarchy}

Hierarchy equations couple:

\[
f_n
\leftrightarrow
f_{n+1}.
\]

Boltzmann equation emerges under molecular chaos assumptions.

\chapter{Hydrodynamics}

\section{Euler Equations}

Fluid density:

\[
\rho(x,t).
\]

Velocity field:

\[
v(x,t).
\]

Continuity equation:

\[
\partial_t\rho
+
\nabla\cdot(\rho v)=0.
\]

Euler equation:

\[
\partial_t v
+
(v\cdot\nabla)v
=
-\frac1\rho\nabla p.
\]

\section{Vorticity}

Vorticity:

\[
\omega=\nabla\times v.
\]

Kelvin circulation theorem follows geometrically.

\part{Operator and Hilbert-Space Formulations}

\chapter{Koopman--von Neumann Mechanics}

\section{Hilbert Space of Classical Mechanics}

Wave functions on phase space:

\[
\psi(q,p).
\]

Inner product:

\[
\langle\phi,\psi\rangle
=
\int
\phi^*\psi\,dqdp.
\]

\section{Liouville Operator}

Dynamics:

\[
i\partial_t\psi
=
\hat L\psi.
\]

where:

\[
\hat L
=
i\{H,\cdot\}.
\]

\section{Importance}

Classical mechanics already possesses:

\begin{itemize}
\item Hilbert spaces,
\item unitary evolution,
\item operator methods.
\end{itemize}

Thus Hilbert structure alone does not imply quantization.

\chapter{Operator and Spectral Methods}

\section{Lax Pairs}

Integrable systems satisfy:

\[
\frac{dL}{dt}=[L,M].
\]

Spectrum of:

\[
L
\]

is conserved.

\section{Inverse Scattering}

Dynamics encoded spectrally.

Examples:

\begin{itemize}
\item KdV equation,
\item nonlinear Schr"odinger equation,
\item Toda lattice.
\end{itemize}

\part{Stochastic and Information-Theoretic Mechanics}

\chapter{Stochastic Mechanics}

\section{Langevin Equation}

Stochastic dynamics:

\[
mdv
=
-Fdt
-\gamma vdt
+\sqrt{2D}\,dW_t.
\]

\section{Fokker--Planck Equation}

Probability density evolves via:

\[
\partial_t\rho
=
-\nabla\cdot(b\rho)
+
D\nabla^2\rho.
\]

\section{Brownian Motion on Manifolds}

Generator becomes Laplace--Beltrami operator:

\[
\Delta_g.
\]

\chapter{Information Geometry and Optimal Transport}

\section{Fisher Geometry}

Probability manifold metric:

\[
g_{ij}
=
\int
\rho
\partial_i\ln\rho
\partial_j\ln\rho
 dx.
\]

\section{Wasserstein Geometry}

Optimal transport metric:

\[
W_2(\rho_1,\rho_2).
\]

Hydrodynamics can be reformulated as geodesic flow on probability space.

\part{Topological and Category-Theoretic Structures}

\chapter{Topological Mechanics}

\section{Configuration-Space Topology}

Topology influences mechanics through:

\begin{itemize}
\item winding numbers,
\item homotopy classes,
\item monodromy,
\item braid groups.
\end{itemize}

\section{Morse Theory}

Critical points of action functional determine topology of trajectory space.

\chapter{Category-Theoretic Mechanics}

\section{Functorial Dynamics}

Mechanical systems can be organized categorically.

Morphisms preserve:

\begin{itemize}
\item symplectic structure,
\item Poisson structure,
\item variational structure.
\end{itemize}

\section{Symplectic Category}

Objects:

\[
(M,\omega).
\]

Morphisms:

canonical relations.

\part{Comparative Structures}

\chapter{Relations Between Formulations}

\section{Transformation Web}

\begin{center}
\begin{tabular}{lll}
Newton & $\to$ Lagrange & variational principle \\
Lagrange & $\to$ Hamilton & Legendre transform \\
Hamilton & $\to$ HJ & generating function \\
Hamilton & $\to$ Liouville & statistical lift \\
Hamilton & $\to$ KvN & Koopman representation \\
Hydrodynamics & $\to$ optimal transport & Wasserstein geometry
\end{tabular}
\end{center}

\section{Primitive Structures}

Different formulations privilege:

\begin{center}
\begin{tabular}{ll}
Newtonian & trajectories \\
Lagrangian & actions \\
Hamiltonian & symplectic flow \\
Poisson & Lie algebras \\
Hydrodynamic & continuum fields \\
Statistical & measures \\
KvN & Hilbert evolution
\end{tabular}
\end{center}

\chapter{Conclusion}

Classical mechanics is not a single theory but an interconnected network of mathematical languages.

Its formulations involve:

\begin{itemize}
\item differential equations,
\item variational principles,
\item symplectic geometry,
\item Poisson algebras,
\item kinetic theory,
\item probability geometry,
\item operator theory,
\item topology,
\item category theory.
\end{itemize}

Many structures usually considered uniquely quantum already appear classically:

\begin{itemize}
\item Hilbert spaces,
\item operator evolution,
\item noncommutative brackets,
\item geometric phases,
\item stochastic dynamics.
\end{itemize}

The modern view of classical mechanics is therefore fundamentally geometric, statistical, and structural rather than merely Newtonian.

\bibliographystyle{plain}

\begin{thebibliography}{99}

\bibitem{Arnold}
V.~I.~Arnold,
\textit{Mathematical Methods of Classical Mechanics}.

\bibitem{MarsdenRatiu}
J.~Marsden and T.~Ratiu,
\textit{Introduction to Mechanics and Symmetry}.

\bibitem{AbrahamMarsden}
R.~Abraham and J.~Marsden,
\textit{Foundations of Mechanics}.

\bibitem{Goldstein}
H.~Goldstein,
\textit{Classical Mechanics}.

\bibitem{JoseSaletan}
J.~Jose and E.~Saletan,
\textit{Classical Dynamics}.

\bibitem{LibermannMarle}
P.~Libermann and C.~Marle,
\textit{Symplectic Geometry and Analytical Mechanics}.

\bibitem{deLeon}
M.~de Le\'on and P.~Rodrigues,
\textit{Methods of Differential Geometry in Analytical Mechanics}.

\bibitem{Morrison}
P.~Morrison,
``Hamiltonian description of fluid and plasma systems.''

\bibitem{Koopman}
B.~Koopman,
``Hamiltonian systems and transformation in Hilbert space.''

\bibitem{vonNeumann}
J.~von Neumann,
``Zur Operatorenmethode in der klassischen Mechanik.''

\end{thebibliography}

\end{document}
